\documentclass[../../../../main.tex]{subfiles}
\begin{document}

\begin{flushright}
\footnotesize\textit{【自動文字起こし・要確認】}
\end{flushright}

{\bf [解]}
$P(1+2t,\, t,\, -2-t)$，$Q(1+s,\, 2-s,\, -3+s)$，$R(1+u,\, -1+2u,\, u)$とおく．

$Q$，$R$が垂足だから，$\overrightarrow{PQ}\perp m$，$\overrightarrow{PR}\perp n$となる．

\begin{align*}
\begin{pmatrix} s-2t \\ 2-s-t \\ -1+s+t \end{pmatrix}\cdot\begin{pmatrix}1\\-1\\1\end{pmatrix}=0
&\iff (s-2t)-(2-s-t)+(-1+s+t)=0\\
&\iff 3s-3=0 \iff s=1，\\[4pt]
\begin{pmatrix} u-2t \\ -1+2u-t \\ 2+u+t \end{pmatrix}\cdot\begin{pmatrix}1\\2\\1\end{pmatrix}=0
&\iff (u-2t)+2(-1+2u-t)+(2+u+t)=0\\
&\iff 6u-3t=0 \iff u=\frac{1}{2}t
\end{align*}

だから，$(s,u)=\left(1,\dfrac{1}{2}t\right)$である．

\begin{align*}
\overline{PQ}^2&=(s-2t)^2+(s+(t-2))^2+(s+(t-1))^2\\
&=3s^2-6s+6t^2-6t+5\\
&=6t^2-6t+2\quad(s=1\text{を代入})，\\[4pt]
\overline{PR}^2&=(u-2t)^2+(2u-(t+1))^2+(u+(t+2))^2\\
&=6u^2-6tu+6t^2+6t+5\\
&=\frac{9}{2}t^2+6t+5\quad\left(u=\frac{1}{2}t\text{を代入}\right)．
\end{align*}

だから，$f(t)=\overline{PQ}^2+\overline{PR}^2$とすると
\begin{align*}
f(t)=\frac{21}{2}t^2+7．
\end{align*}

これは$t=0$で最小値$7$をとる．この時$P(1,0,-2)$である．

\newpage

\end{document}
